\chapter{Mechanical Design}\label{ch:mechanical-design}

\section{Mechanical Concept}\label{mechanical-concept}

Zephyr is organized about a common horizontal wheel axis. Four
independently driven Mecanum wheels are distributed along this line, and
the body mass is carried above the line so that gravity produces an
unstable pitch moment. In the nominal upright configuration,
longitudinal motion is perpendicular to the common wheel axis, lateral
motion is parallel to that axis, yaw is rotation about the vertical
axis, and pitch is rotation about the common wheel axis. This definition
must be reconciled with the final CAD coordinate system and IMU mounting
axes before signs are frozen.

A rectangular four-corner Mecanum chassis obtains yaw authority from
longitudinal and lateral offsets. In the collinear design, all
wheel-center longitudinal coordinates are nominally zero. Yaw authority
must therefore come from the differences in longitudinal force applied
at distinct axial positions. Lateral force depends on the roller
handedness. The physical order of left- and right-handed wheels is
consequently a functional design variable rather than a cosmetic
assembly detail.

\begin{longtable}[]{@{}
  >{\raggedright\arraybackslash}p{(\columnwidth - 0\tabcolsep) * \real{1.0000}}@{}}
\toprule\noalign{}
\begin{minipage}[b]{\linewidth}\raggedright
\textless INSERT FIGURE 4.1: Insert a high-resolution photograph of the
current mechanical prototype on a plain background. Label wheel numbers,
actuator bodies, structural frame, battery location, electronics
enclosure, IMU location, cable paths, and balance axis. Identify which
components are temporary\textgreater{}\\
What the figure should show: Insert a high-resolution photograph of the
current mechanical prototype on a plain background. Label wheel numbers,
actuator bodies, structural frame, battery location, electronics
enclosure, IMU location, cable paths, and balance axis. Identify which
components are temporary\\
Likely source: physical prototype photograph\strut
\end{minipage} \\
\midrule\noalign{}
\endhead
\bottomrule\noalign{}
\endlastfoot
\end{longtable}

\begin{center}\singlespacing
\phantomsection\addcontentsline{lof}{figure}{\protect\numberline{4.1}Current Zephyr mechanical prototype with major components identified.}
\textbf{Figure 4.1. Current Zephyr mechanical prototype with major components identified.}
\end{center}

\begin{longtable}[]{@{}
  >{\raggedright\arraybackslash}p{(\columnwidth - 0\tabcolsep) * \real{1.0000}}@{}}
\toprule\noalign{}
\begin{minipage}[b]{\linewidth}\raggedright
\textless INSERT FIGURE 4.2: Insert an exploded CAD view with a
bill-of-material callout for frame plates, printed parts, actuator
mounts, wheel hubs, bearings or supports, fasteners, battery retention,
electronics mounts, covers, and cable restraints\textgreater{}\\
What the figure should show: Insert an exploded CAD view with a
bill-of-material callout for frame plates, printed parts, actuator
mounts, wheel hubs, bearings or supports, fasteners, battery retention,
electronics mounts, covers, and cable restraints\\
Likely source: final SolidWorks assembly and drawing package\strut
\end{minipage} \\
\midrule\noalign{}
\endhead
\bottomrule\noalign{}
\endlastfoot
\end{longtable}

\begin{center}\singlespacing
\phantomsection\addcontentsline{lof}{figure}{\protect\numberline{4.2}Exploded CAD view of the chassis, actuator mounts, wheels, and electronics structure.}
\textbf{Figure 4.2. Exploded CAD view of the chassis, actuator mounts, wheels, and electronics structure.}
\end{center}

\section{Coordinate Frames and Wheel Arrangement}\label{coordinate-frames-and-wheel-arrangement}

The body frame \{B\} is placed on the common wheel-center line at a
defined axial origin. The x\_B axis points in the nominal forward
direction, y\_B follows the common wheel axis, and z\_B points upward
when the robot is upright. Wheel i is located at p\_i=[0, y\_i, 0]\ensuremath{^{T}}
for the ideal collinear model. A roller sign s\_i is assigned from the
direction in which positive wheel rotation couples to positive lateral
velocity. The sign must be determined by inspection and verified by
lifting the platform and commanding one wheel at a time.

The balance angle \ensuremath{\theta} is zero when the designated body reference plane is
vertical. Positive pitch, yaw, wheel rotation, motor current, and torque
directions must be documented on the same drawing. The IMU mounting
rotation is then a fixed transform from the sensor frame to \{B\};
software should apply this transform before the estimator. Using one
signed convention throughout CAD, wiring, firmware, simulation, and
plots prevents the common error of compensating for a reversed motor in
multiple layers.

\begin{longtable}[]{@{}
  >{\raggedright\arraybackslash}p{(\columnwidth - 0\tabcolsep) * \real{1.0000}}@{}}
\toprule\noalign{}
\begin{minipage}[b]{\linewidth}\raggedright
\textless INSERT FIGURE 4.3: Provide a dimensioned coordinate-frame
drawing with the origin, x\_B-y\_B-z\_B directions, inertial frame,
positive pitch and yaw, wheel centers y1 through y4, effective wheel
radius, roller angle and sign, center-of-mass coordinates, and ground
plane\textgreater{}\\
What the figure should show: Provide a dimensioned coordinate-frame
drawing with the origin, x\_B-y\_B-z\_B directions, inertial frame,
positive pitch and yaw, wheel centers y1 through y4, effective wheel
radius, roller angle and sign, center-of-mass coordinates, and ground
plane\\
Likely source: CAD drawing supplemented by MATLAB vector graphics\strut
\end{minipage} \\
\midrule\noalign{}
\endhead
\bottomrule\noalign{}
\endlastfoot
\end{longtable}

\begin{center}\singlespacing
\phantomsection\addcontentsline{lof}{figure}{\protect\numberline{4.3}Body coordinate frame, collinear wheel locations, roller signs, and balancing axis.}
\textbf{Figure 4.3. Body coordinate frame, collinear wheel locations, roller signs, and balancing axis.}
\end{center}

\begin{longtable}[]{@{}
  >{\raggedright\arraybackslash}p{(\columnwidth - 0\tabcolsep) * \real{1.0000}}@{}}
\toprule\noalign{}
\begin{minipage}[b]{\linewidth}\raggedright
\textless AUTHOR INPUT REQUIRED: Explain why the collinear four-wheel
arrangement was selected instead of a conventional two-wheel balancing
platform or a statically stable four-corner Mecanum platform. Address
expected lateral authority, packaging, redundancy, research value, and
accepted mechanical complexity.\textgreater{}
\end{minipage} \\
\midrule\noalign{}
\endhead
\bottomrule\noalign{}
\endlastfoot
\end{longtable}

\section{Wheels and Actuators}\label{wheels-and-actuators}

The selected wheels are the goBILDA GripForce Mecanum Wheel Set, whose
manufacturer lists a 104 mm diameter, 40A-durometer silicone rollers, 11
bearing-supported rollers per wheel, and a mass of 236 g per wheel
\cite{gobilda2026gripforce}. The 52 mm nominal radius inferred from the catalog diameter is
suitable for preliminary calculations, but the effective rolling radius
under robot load should be measured because roller deformation and floor
compliance change the kinematic scale.

Each wheel is driven by a CubeMars AK40-10 actuator. The compact
integrated package reduces the number of separate motor-controller and
encoder mounts. It also imposes a 10:1 planetary gearbox and a single
internal encoder in the documented variant \cite{cubemars2026ak4010}. The manufacturer
lists 18 arcmin of backlash. At the wheel rim, the corresponding lost
motion depends on the actual output backlash and wheel radius; this
value should be measured rather than inferred from the catalog alone.
Backlash is particularly important near the zero-speed reversals
generated by a balance controller.

\begin{longtable}[]{@{}
  >{\raggedright\arraybackslash}p{(\columnwidth - 0\tabcolsep) * \real{1.0000}}@{}}
\toprule\noalign{}
\begin{minipage}[b]{\linewidth}\raggedright
\textless INSERT FIGURE 4.4: Show a section or exploded detail from
actuator housing to wheel: actuator mounting face, output interface, hub
or adapter, fasteners, wheel, any shaft support or bearing, axial
retention, locating features, and torque reaction into the frame.
Include critical dimensions and materials\textgreater{}\\
What the figure should show: Show a section or exploded detail from
actuator housing to wheel: actuator mounting face, output interface, hub
or adapter, fasteners, wheel, any shaft support or bearing, axial
retention, locating features, and torque reaction into the frame.
Include critical dimensions and materials\\
Likely source: manufacturing drawing and actuator\slash{}wheel CAD\strut
\end{minipage} \\
\midrule\noalign{}
\endhead
\bottomrule\noalign{}
\endlastfoot
\end{longtable}

\begin{center}\singlespacing
\phantomsection\addcontentsline{lof}{figure}{\protect\numberline{4.4}Actuator-to-wheel structural interface and load path.}
\textbf{Figure 4.4. Actuator-to-wheel structural interface and load path.}
\end{center}

\section{Chassis, Structural Interfaces, and Mass Distribution}\label{chassis-structural-interfaces-and-mass-distribution}

The chassis must carry the battery and electronics rigidly relative to
the IMU while transmitting wheel reaction torque among four actuator
mounts. Pitch-control bandwidth can excite compliance if the body twists
about the common axis or if the IMU mount flexes. The structural design
should therefore use closed load paths, short actuator-mount spans,
positive locating features, and fasteners sized for repeated reversals.
Removable electronics and battery modules improve serviceability, but
their interfaces must not introduce motion into the IMU reference.

Center-of-mass placement creates a deliberate tradeoff. A higher center
of mass increases the gravitational moment for a given pitch angle and
generally slows the open-loop falling motion, but it also increases the
torque and energy required for recovery. A lower center of mass reduces
gravitational leverage but can demand faster sensing and actuation. The
final controller model requires the assembled mass, center-of-mass
height and offsets, pitch inertia, and wheel\slash{}actuator inertias. CAD mass
properties should be checked by physical weighing and a center-of-mass
or pendulum test because fasteners, wiring, printed infill, and
batteries are often modeled inaccurately.

\clearpage
\begin{landscape}
\par\medskip\noindent\singlespacing
\phantomsection\addcontentsline{lot}{table}{\protect\numberline{4.1}Mechanical parameters required for analysis and validation}
\textbf{Table 4.1. Mechanical parameters required for analysis and validation}\par

\begin{longtable}[]{@{}
  >{\raggedright\arraybackslash}p{(\columnwidth - 8\tabcolsep) * \real{0.2266}}
  >{\raggedright\arraybackslash}p{(\columnwidth - 8\tabcolsep) * \real{0.1250}}
  >{\raggedright\arraybackslash}p{(\columnwidth - 8\tabcolsep) * \real{0.1094}}
  >{\raggedright\arraybackslash}p{(\columnwidth - 8\tabcolsep) * \real{0.3359}}
  >{\raggedright\arraybackslash}p{(\columnwidth - 8\tabcolsep) * \real{0.2031}}@{}}
\toprule\noalign{}
\begin{minipage}[b]{\linewidth}\raggedright
\textbf{Parameter}
\end{minipage} & \begin{minipage}[b]{\linewidth}\raggedright
\textbf{Symbol}
\end{minipage} & \begin{minipage}[b]{\linewidth}\raggedright
\textbf{Unit}
\end{minipage} & \begin{minipage}[b]{\linewidth}\raggedright
\textbf{Required method}
\end{minipage} & \begin{minipage}[b]{\linewidth}\raggedright
\textbf{Current value\slash{}status}
\end{minipage} \\
\midrule\noalign{}
\endhead
\bottomrule\noalign{}
\endlastfoot
Assembled robot mass & m & kg & Calibrated scale with final battery and
wiring & \textless MEASURE\textgreater{} \\
Body mass above wheel axis & m\_b & kg & CAD grouping and component
weighing & \textless MEASURE\textgreater{} \\
Effective wheel radius & r & m & Loaded rollout distance over multiple
revolutions & Catalog radius 0.052 m; effective value unverified \\
Wheel axial locations & y1...y4 & m & CAD and physical measurement from
the defined origin & \textless MEASURE\textgreater{} \\
Center-of-mass height & \ensuremath{\ell} & m & CAD mass properties plus physical
confirmation & \textless MEASURE\textgreater{} \\
COM x and y offsets & x\_c, y\_c & m & Balance\slash{}knife-edge or suspension
method & \textless MEASURE\textgreater{} \\
Pitch moment of inertia & I\_p & kg\ensuremath{\cdot}m\ensuremath{^{2}} & CAD plus bifilar\slash{}trifilar or
system-identification test & \textless MEASURE\textgreater{} \\
Yaw moment of inertia & I\_z & kg\ensuremath{\cdot}m\ensuremath{^{2}} & CAD plus torsional test or system
identification & \textless MEASURE\textgreater{} \\
Wheel\slash{}rotor equivalent inertia & J\_eq & kg\ensuremath{\cdot}m\ensuremath{^{2}} & Manufacturer data and
reflected-inertia calculation & \textless VERIFY\textgreater{} \\
Roller angle and signs & \ensuremath{\beta}, s\_i & deg, --- & Physical inspection and
single-wheel direction test & \textless VERIFY FROM
HARDWARE\textgreater{} \\
\end{longtable}
\end{landscape}
\clearpage

\section{Manufacturability, Assembly, and Serviceability}\label{manufacturability-assembly-and-serviceability}

The final drawing package should control actuator location, wheel
spacing, perpendicularity between the wheel axis and body reference, and
the orientation of each Mecanum wheel. Assembly should use an
unambiguous wheel-number label and roller-direction mark that remains
visible after installation. Threaded joints subject to reversing torque
should use appropriate locking features and documented tightening
practices. The battery should be mechanically retained independently of
its electrical connector, and cables should include strain relief and
clearance from rotating rollers.

Service access should allow the battery, controller, CAN transceiver,
receiver, display, and IMU to be removed without disturbing the entire
drivetrain. Connectors should be keyed or labeled so that a CAN plug
cannot be inserted into a power header. Test points for logic rails,
CANH, CANL, and ground should remain accessible when the robot is
restrained. A removable physical stand or tether attachment is
recommended for initial controller tests.

\section{Mechanical Status and Remaining Evidence}\label{mechanical-status-and-remaining-evidence}

\par\medskip\noindent\singlespacing
\phantomsection\addcontentsline{lot}{table}{\protect\numberline{4.2}Mechanical component and interface status}
\textbf{Table 4.2. Mechanical component and interface status}\par

\begin{longtable}[]{@{}
  >{\raggedright\arraybackslash}p{(\columnwidth - 6\tabcolsep) * \real{0.2000}}
  >{\raggedright\arraybackslash}p{(\columnwidth - 6\tabcolsep) * \real{0.3200}}
  >{\raggedright\arraybackslash}p{(\columnwidth - 6\tabcolsep) * \real{0.2000}}
  >{\raggedright\arraybackslash}p{(\columnwidth - 6\tabcolsep) * \real{0.2800}}@{}}
\toprule\noalign{}
\begin{minipage}[b]{\linewidth}\raggedright
\textbf{Item}
\end{minipage} & \begin{minipage}[b]{\linewidth}\raggedright
\textbf{Evidence available}
\end{minipage} & \begin{minipage}[b]{\linewidth}\raggedright
\textbf{Apparent status}
\end{minipage} & \begin{minipage}[b]{\linewidth}\raggedright
\textbf{Information required}
\end{minipage} \\
\midrule\noalign{}
\endhead
\bottomrule\noalign{}
\endlastfoot
Four-wheel collinear concept & Project description and author statements
& Selected & Dimensioned drawing and roller order \\
AK40-10 actuators & Repository firmware and manufacturer documentation &
Motors reported moving & Installed revision, mount drawing, output
interface \\
104 mm Mecanum wheels & Selected part and manufacturer specification &
Selected\slash{}procured & Installed handedness and loaded radius \\
Frame\slash{}chassis & Author reports mechanical build approximately 80 percent
complete & Partially manufactured & CAD, materials, dimensions, mass
properties, photos \\
Bearings\slash{}shafts\slash{}couplings & No sufficient repository evidence available
& Unconfirmed & Part numbers, fits, retention, calculations \\
Cable management and covers & No sufficient final evidence available &
Pending or unconfirmed & Photos and assembly procedure \\
Structural analysis & No confirmed FEA or hand calculations & Pending &
Load cases, material properties, stresses, deflections, safety
factors \\
\end{longtable}
